\section{Python Simulation}%
\label{sec:simulation}
\emph{A python simulation has been created that will simulate the gas flow of the hydrogen power system, which provides insight to the behaviour of this complex system. This section describes the goal of the simulation, how it works and the results it produced.}

\subsection*{Goal}
The goal of the SHARP-PS project is to develop a power system for a UAV that produces electricity from hydrogen using a fuel cell. This hydrogen will not be supplied by a hydrogen tank, but by a Cool Gas Generator (CGG) which will be developed by Solid Flow.

This system is quite complex and has many open variables that all influence each other.
For example, both the hydrogen generation and consumption rate affect the pressure of the hydrogen supply system. The hydrogen consumption rate is in turn affected by the frequency and length of purge events and the power demand on the fuel cell. The hydrogen production rate is dependent on the size of the CGG, but the size of the CGG also influences the volume available for the hydrogen thus also affecting the pressure in a different way.

In order to grasp this complex system of variables, a python simulation has been created that simulates the gas flows of the entire system. This allows the effect of changing each variable to be observed, which will be a useful tool during the development of the system.

\subsection*{Process}
The core of the simulation are a set of volumes that contain either hydrogen or electricity, and several processes that move content from one volume to another at a rate determined by the system properties. Figure~\ref{fig:sim diagram} provides an overview of these volumes and processes, and the parameters each depends on.

\newcommand{\pytable}[2]{
    \begin{tabular}{c}
        \textbf{#1} \\
        \midrule
        #2          \\
    \end{tabular}
}

\newcommand{\cggstorage}{
    \pytable{CGG Storage}{Capacity \\ Amount}
}

\newcommand{\cggproc}{
    \pytable{CGG Production}{Rate}
}

\newcommand{\buffer}{
    \pytable{Buffer Volume}{Volume}
}

\newcommand{\fuelcell}{
    \pytable{Fuel Cell}{max power \\ min pressure \\ max pressure \\ normal massflow \\ purge massflow \\ purge duration \\ purge interval}
}

\newcommand{\control}{
    \pytable{Controller}{control scheme}
}

\newcommand{\battery}{
    \pytable{Battery}{Capacity}
}

\newcommand{\drone}{
    \pytable{Drone}{Climb Power\\ Cruise Power \\ Flight Profile}
}

\newcommand{\relief}{
    \pytable{Relief Valve}{Relief Pressure}
}

\begin{figure}[htp]
    \centering
    \scalebox{0.6}{
        \begin{tikzpicture}[node distance=37mm, anchor=north]
            \node (v_cgg) [pyvolume] {\cggstorage};
            \node (p_cgg) [pyprocess, right of=v_cgg] {\cggproc};
            \node (v_buffer) [pyvolume, right of=p_cgg] {\buffer};
            \node (p_fc) [pyprocess, right of=v_buffer] {\fuelcell};
            \node (p_control) [pyprocess, right of=p_fc] {\control};
            \node (v_battery) [pyvolume, right of=p_control] {\battery};
            \node (p_drone) [pyprocess, right of=v_battery] {\drone};

            \node (p_relief) [pyprocess, below of=v_buffer, yshift=17mm] {\relief};
            \node (v_relief) [pyvolume, below of=p_relief, yshift=2cm] {\textbf{Wasted H2}};

            \draw [arrow] (v_cgg) -- (p_cgg);
            \draw [arrow] (p_cgg) -- (v_buffer);
            \draw [arrow] (v_buffer) -- (p_fc);
            \draw [arrow] (p_fc) -- (p_control);
            \draw [arrow] (p_control) -- (v_battery);
            \draw [arrow] (v_battery) -- (p_drone);

            \draw [arrow] (v_buffer) -- (p_relief);
            \draw [arrow] (p_relief) -- (v_relief);

            \draw [arrow] (v_buffer) -- ++(0,3) -| ($(p_control.north) - (0.2,0)$);
            \draw [arrow] (v_battery) -- ++(0,3) -| ($(p_control.north) + (0.2,0)$);
        \end{tikzpicture}
    }
    \caption{Simulation Diagram}
    \label{fig:sim diagram}
\end{figure}

\subsection*{Implementation}
The initial parameters chosen for the simulation are based on the mission scenario by Wren Aerospace. These parameters are listed in table~\ref{tab:direct_sim_parameters}. From these parameters, a preliminary fuel cell can be selected which will provide typical parameters for fuel cells in this range, visible in table~\ref{tab:derived_sim_parameters}. The final selecting of the fuel cell to be used will be discussed later in section~\ref{sec:h2 design}. The simulation parameters may be adjusted at any time to suit the selected fuel cell.

\newcommand{\simpar}[3]{#1 & #2 & #3\\}

\begin{table}[htp]
    \centering
    \begin{subtable}[t]{0.3\textwidth}
        \centering
        \scalebox{0.85}{
            \begin{tabular}[t]{lrl}
                \toprule
                \textbf{Parameter} & \textbf{Value} & \textbf{Unit} \\
                \midrule
                \simpar{CGG Capacity}{160}{g}
                \simpar{Fuel Cell max Power}{200}{W}
                \simpar{Climb Power}{200--220}{W}
                \simpar{Cruise}{140--180}{W}
                \simpar{Flight time}{15}{h}
                \bottomrule
            \end{tabular}}
        \caption{Based on the Mission Scenario by Wren Aerospace}%
        \label{tab:direct_sim_parameters}%
    \end{subtable}%
    \hspace{\fill}
    \begin{subtable}[t]{0.3\textwidth}
        \centering
        \scalebox{0.85}{
            \begin{tabular}[t]{lrl}
                \toprule
                \textbf{Parameter} & \textbf{Value} & \textbf{Unit} \\
                \midrule
                \simpar{Min pressure}{0.35}{Bar}
                \simpar{Max pressure}{0.5}{Bar}
                \simpar{Normal massflow}{0.263}{g/min}
                \simpar{Purge massflow}{0.545}{g/min}
                \simpar{Purge duration}{1--2}{s}
                \simpar{Purge interval}{30--90}{s}
                \bottomrule
            \end{tabular}}%
        \caption{Derived from the selected Fuel Cell}%
        \label{tab:derived_sim_parameters}%
    \end{subtable}
    \hspace{\fill}
    \begin{subtable}[t]{0.3\textwidth}
        \centering
        % \flushright
        \scalebox{0.85}{
            \begin{tabular}[t]{lrl}
                \toprule
                \textbf{Parameter} & \textbf{Value} & \textbf{Unit} \\
                \midrule
                \simpar{CGG Amount}{1}{---}
                \simpar{CGG rate}{0.21}{g/min}
                \simpar{Buffer Volume}{0.1}{L}
                \simpar{Relief Pressure}{4}{Bar}
                \simpar{Battery Capacity}{500}{kJ}
                \bottomrule
            \end{tabular}}
        \caption{Chosen}%
        \label{tab:chosen_sim_parameters}%
    \end{subtable}%
    \caption{Simulation Parameters}%
    \label{tab:sim_parameters}%
\end{table}

% \emph{Something about control scheme and flight profile...}
The flight profile has not been determined yet, for the purposes of this simulation a climb of 1 hour has been chosen after which the drone will cruise for the rest of the mission.

The purpose of the controller is to limit the hydrogen consumption when the pressure in the buffer volume gets too low or when the battery is fully charged. For this purpose both an On-Off controller and a PID controller have been tested.

% \todo{Discuss controller}
\newpage
